% Generated draft: requires mathematical and source review.
\documentclass[11pt]{article}
\usepackage[margin=25mm]{geometry}
\usepackage{fontspec}
\setmainfont{Latin Modern Roman}
\usepackage{amsmath,amssymb,mathtools,unicode-math}
\setmathfont{Latin Modern Math}
\usepackage{enumitem,needspace,xurl,hyperref}
\setlength{\parindent}{0pt}
\setlength{\parskip}{5pt}
\setlength{\emergencystretch}{3em}
\newcommand{\N}{\mathbb{N}}
\newcommand{\Z}{\mathbb{Z}}
\newcommand{\Q}{\mathbb{Q}}
\newcommand{\R}{\mathbb{R}}
\newcommand{\C}{\mathbb{C}}
\newcommand{\F}{\mathbb{F}}
\newcommand{\A}{\mathbb{A}}
\newcommand{\PP}{\mathbb{P}}
\newcommand{\E}{\mathbb{E}}
\newcommand{\eps}{\varepsilon}
\newcommand{\dd}{\,\mathrm{d}}
\begin{document}
\section*{UnsolvedMath ID: 11100040}
\textbf{Source problem code:} catalogue entry
\subsection{Definitions and mathematical statement}
A Grothendieck abelian category \(\mathcal A\) has all small colimits, exact filtered colimits and a generator. The source asks for an axiomatic understanding of generalized homology theories valued in such categories, motivated by Grojnowski's equivariant elliptic cohomology, which it describes tentatively as sheaf-valued rather than module-valued over one ring.

A basic target-dependent formulation would specify functors \(h_n\) from the chosen spaces/pairs to \(\mathcal A\), with natural boundary morphisms, homotopy invariance, exact sequences for pairs, excision and the appropriate additivity/suspension axioms. For a cohomology theory one instead uses contravariant functors and the corresponding cohomological exactness. The source does not choose these axioms, the equivariant input category, graded versus derived values, representability notion or tensor/localization structure. Merely replacing abelian groups by \(\mathcal A\) in an exact sequence does not by itself identify how the motivating elliptic theory is represented or characterized.

The intended framework should explain Grojnowski's example with a specified sheaf category and base parameter space, and clarify which familiar homology-theory axioms and structural results remain appropriate. A coherent-sheaf category may need embedding into a larger ambient sheaf category before Grothendieck-category axioms apply; those are not interchangeable automatically. The source's ``scheme or something'' is an expressed uncertainty, not a fixed choice between algebraic and analytic sheaves. Likewise a sheaf category need not be identified with a single ring's module category by simply taking global sections. No complete categorical equivalence or universal characterization is asserted as the requested answer.
\subsection{Short English statement}
Develop suitable axioms and structural tools for homology or cohomology theories valued in Grothendieck abelian categories, explaining the sheaf-valued equivariant elliptic example. The precise target, input category and representation framework remain unspecified.
\subsection{Sources}
\begin{itemize}
\item Local UnsolvedMath registry, permanent ID \(11100040\), exact source code \texttt{catalogue entry}.
\item Mark Hovey, \emph{Hovey's Algebraic Topology Problem List}, Axiomatic stable homotopy, item 8, surviving mirror: \url{https://www-users.cse.umn.edu/~tlawson/hovey/}. Registry access 2026-07-29; html-ordered-list extraction.
\item Grojnowski's equivariant elliptic cohomology is the motivating attribution, but no exact paper or target-space/sheaf definition is supplied.
\end{itemize}
\end{document}
